\documentclass[11pt]{article}
\usepackage[margin=1in]{geometry}
\usepackage{amsmath,amssymb,amsthm}
\usepackage{booktabs}
\usepackage{xcolor}
\usepackage[colorlinks=true,linkcolor=blue!50!black,urlcolor=blue!50!black,citecolor=blue!50!black]{hyperref}
\usepackage{url}
\usepackage{listings}
\emergencystretch=3em
\lstset{basicstyle=\ttfamily\small,columns=fullflexible,keepspaces=true,frame=single,
  framerule=0.3pt,xleftmargin=6pt,xrightmargin=6pt,literate={→}{{$\to$}}1 {↔}{{$\leftrightarrow$}}1
  {ℕ}{{$\mathbb{N}$}}1 {∃}{{$\exists$}}1 {∀}{{$\forall$}}1 {∧}{{$\wedge$}}1 {∨}{{$\vee$}}1}

\newtheorem{theorem}{Theorem}
\newtheorem{lemma}[theorem]{Lemma}
\newtheorem{corollary}[theorem]{Corollary}
\theoremstyle{remark}
\newtheorem{remark}[theorem]{Remark}

\newcommand{\ord}{\operatorname{ord}}
\newcommand{\lcm}{\operatorname{lcm}}

\title{Consecutive residues of equal multiplicative order:\\
a proof of the conjecture in OEIS A397683}
\author{Vasily Shablov\\[2pt]
\small Proof search and Lean formalization assisted by Claude (Anthropic)}
\date{October 2026}

\begin{document}
\maketitle

\begin{abstract}
For an integer $m\ge 2$ let $f(m)$ be the number of $x$ with $1\le x<m$ such that $x$ and
$x+1$ are both units modulo $m$ and $\ord_m(x)=\ord_m(x+1)$. OEIS sequence A397683
(D.~S.~Gonz\'alez May, September 2026) conjectures that for odd $m>1$, $f(m)=0$ if and only if
$m$ is a power of $3$ or a power of $7$. We prove this conjecture. The only non-elementary
input is S.~D.~Cohen's theorem that every prime $p>7$ has a pair of consecutive primitive
roots. The rest is an explicit Chinese remainder construction, with the residues $4$ modulo
$3^a$ and $3$ modulo $7^b$ chosen so that the multiplicative orders balance. The full argument
is formalized in Lean~4 with Mathlib. Cohen's theorem enters as an explicit hypothesis and
nothing else is assumed.
\end{abstract}

\section{Introduction}

For $m\ge 2$ define
\[
  f(m)=\#\{\,1\le x<m:\ \gcd(x,m)=\gcd(x+1,m)=1,\ \ord_m(x)=\ord_m(x+1)\,\}.
\]
Every unit modulo an even $m$ is odd, so $x$ and $x+1$ cannot both be units and $f(m)=0$.
The interesting case is odd $m$. OEIS A397683 lists $a(n)=f(2n-1)$ and records the following.

\paragraph{Conjecture (A397683).} For odd $m$, $f(m)=0$ if and only if $m$ is a power of $3$
or a power of $7$. This was verified for $m<20000$.

Here ``power'' means $3^a$ or $7^b$ with exponent at least $1$, since $f(1)=1$ by convention.
The entry already notes that positivity is preserved under coprime products: if
$\gcd(r,s)=1$, $f(r)>0$ and $f(s)>0$, then $f(rs)>0$. Together with Cohen's theorem this settles
every odd $m$ coprime to $21$. The open part is every odd $m$ divisible by $3$ or $7$ that is
not a pure power of $3$ or of $7$. There the coprime-product argument fails, because
$f(3^a)=f(7^b)=0$. We close this case.

\begin{theorem}\label{thm:main}
Assume Cohen's theorem (Theorem~\ref{thm:cohen}). For every odd $m>1$, $f(m)=0$ if and only if
$m=3^a$ or $m=7^b$ for some $a,b\ge 1$.
\end{theorem}

\paragraph{Prior work.} We checked the following on 9 October 2026 and found no existing proof:
\begin{itemize}
  \item The OEIS entry and its edit history. The last edit was on 20 September 2026, and the
    conjecture is still marked open. No other OEIS entry references A397683.
  \item Google DeepMind's \texttt{formal-conjectures} repository, including its issues and pull
    requests.
  \item The \texttt{the-omega-institute/trureturing} repository, an automated pipeline that
    proves OEIS conjectures, including its issues and pull requests.
  \item Recent papers proving OEIS conjectures, namely the ``OEIS Open'' benchmark and S.~Fried's
    series. Both predate the entry.
  \item General web searches.
\end{itemize}
N.~A.~Carella's paper~\cite{carella}, which the entry cites, concerns consecutive primitive
roots modulo primes only.

\section{Tools}

\begin{theorem}[Cohen~\cite{cohen1,cohen2}; see also~\cite{cot}]\label{thm:cohen}
For every prime power $q>7$, the field $\mathbb F_q$ contains two consecutive primitive
elements. In particular, every prime $p>7$ has consecutive primitive roots $y,\,y+1$ modulo
$p$.
\end{theorem}

For $p=5$ the pair $(2,3)$ works, because both residues have order $4$. We write $E_n(z,k)$
for $z^k\equiv 1\pmod n$, so that $\ord_n(z)\mid k\iff E_n(z,k)$. In particular,
$\ord_m(x)=\ord_m(x+1)$ holds if and only if $E_m(x,k)\iff E_m(x+1,k)$ for every $k$.

\begin{lemma}[Order profile]\label{lem:profile}
Let $p$ be an odd prime and $p\nmid z$. Let $d$ be the order of $z$ modulo $p$, and suppose
$p^2\nmid z^d-1$. Then for every $n\ge 1$ and every $k$,
\[ z^k\equiv 1 \pmod{p^n}\iff d\,p^{\,n-1}\mid k. \]
\end{lemma}
\begin{proof}
Both sides force $d\mid k$. Write $k=dj$ with $j\ge1$; the case $j=0$ is trivial. By the
lifting-the-exponent lemma,
$v_p\bigl((z^d)^j-1\bigr)=v_p(z^d-1)+v_p(j)=1+v_p(j)$. Hence $p^n\mid z^k-1$ if and only if
$p^{n-1}\mid j$, which holds if and only if $dp^{n-1}\mid k$.
\end{proof}

\begin{lemma}[Lifting a consecutive pair]\label{lem:lift}
Let $p\ge5$ be prime, and let $y,\,y+1$ be primitive roots modulo $p$. Then there is
$t\in\{0,1,2\}$ such that $y'=y+tp$ and $y'+1$ both satisfy $p^2\nmid (\cdot)^{p-1}-1$.
Consequently both $y'$ and $y'+1$ have order $\varphi(p^n)=(p-1)p^{n-1}$ modulo $p^n$ for every
$n\ge1$.
\end{lemma}
\begin{proof}
Modulo $p^2$, $(y+tp)^{p-1}\equiv y^{p-1}+(p-1)y^{p-2}tp$. Suppose two values $t_1\ne t_2$
both give $(y+t_ip)^{p-1}\equiv1\pmod{p^2}$. Then $p\mid (p-1)y^{p-2}(t_1-t_2)$, so
$p\mid t_1-t_2$. This is impossible, since $|t_1-t_2|\le2<p$. So at most one $t$ is bad for
$y$, and likewise at most one for $y+1$, which leaves a good $t$ among the three. The order
statement then follows from Lemma~\ref{lem:profile} with $d=p-1$.
\end{proof}

\begin{lemma}[Profiles modulo $3^a$ and $7^b$]\label{lem:37}
For $a,b\ge1$ and all $k$:
\begin{align*}
4^k\equiv1 \!\!\pmod{3^a} &\iff 3^{a-1}\mid k, &
5^k\equiv1 \!\!\pmod{3^a} &\iff 2\cdot3^{a-1}\mid k,\\
3^k\equiv1 \!\!\pmod{7^b} &\iff 6\cdot7^{b-1}\mid k, &
4^k\equiv1 \!\!\pmod{7^b} &\iff 3\cdot7^{b-1}\mid k.
\end{align*}
\end{lemma}
\begin{proof}
Apply Lemma~\ref{lem:profile}. The orders modulo the prime are
$\ord_3(4)=1$, $\ord_3(5)=2$, $\ord_7(3)=6$ and $\ord_7(4)=3$. The non-divisibility conditions
are $9\nmid 3$, $9\nmid 24$, $49\nmid 728$ and $49\nmid 63$.
\end{proof}

\section{Proof of Theorem~\ref{thm:main}}

\paragraph{Pure powers have $f=0$.}
\begin{itemize}
  \item \emph{$m=3^a$, $a\ge1$.} If $x$ and $x+1$ are units modulo $3$, then $x\equiv1\pmod 3$.
    So $x^{3^{a-1}}\equiv1\pmod{3^a}$. However, $3^{a-1}$ is odd, so
    $(x+1)^{3^{a-1}}\equiv2^{3^{a-1}}\equiv2\pmod 3$, and $(x+1)^{3^{a-1}}\not\equiv1$.
    The exponent $k=3^{a-1}$ separates the two orders.
  \item \emph{$m=7^b$, $b\ge1$.} Here $x\bmod 7\in\{1,\dots,5\}$. The orders modulo $7$ of the
    consecutive pairs $(x,x+1)$ are
    \[ (1,3),\ (3,6),\ (6,3),\ (3,6),\ (6,2). \]
    In each case there is a $d\in\{1,3,2\}$ such that exactly one of $x^d$, $(x+1)^d$ is
    $\equiv1\pmod7$. Take $k=d\cdot7^{b-1}$. The side with $z^d\equiv1\pmod 7$ satisfies
    $z^k\equiv1\pmod{7^b}$. On the other side, $7^{b-1}\equiv1\pmod 6$ gives
    $w^k\equiv w^d\not\equiv1\pmod 7$. Again the orders differ.
\end{itemize}

\paragraph{Every other odd $m>1$ has $f(m)>0$.}
Write $m=3^a\,7^b\,r$ with $\gcd(r,21)=1$. Since $m$ is not a pure power, either $a,b\ge1$ or
$r>1$.

\emph{The part $r$.} For each $p^n\,\|\,r$ we have $p\ge5$ and $p\ne7$. Theorem~\ref{thm:cohen}
(or the pair $(2,3)$ when $p=5$) together with Lemma~\ref{lem:lift} gives $y_p$ with
$y_p,\,y_p+1$ both of order $\varphi(p^n)$ modulo $p^n$. Glue these by the CRT into $y_r$. Then
for every $k$, $E_r(y_r,k)\iff E_r(y_r+1,k)$, and $E_r(y_r,k)$ forces $2\mid k$ when $r>1$,
because every $\varphi(p^n)$ is even.

\emph{The construction.} By the CRT choose $x$ with
\[ x\equiv4\ (3^a),\qquad x\equiv3\ (7^b),\qquad x\equiv y_r\ (r). \]
Then $x+1\equiv5,\ 4,\ y_r+1$ respectively, so $x$ and $x+1$ are units modulo $m$. By
Lemma~\ref{lem:37}, writing $R(k)$ for $E_r(y_r,k)$:
\begin{align*}
E_m(x,k)&\iff [a=0 \vee 3^{a-1}\mid k]\wedge[b=0\vee (2\mid k\wedge 3\cdot7^{b-1}\mid k)]
  \wedge R(k),\\
E_m(x+1,k)&\iff [a=0 \vee (2\mid k\wedge 3^{a-1}\mid k)]\wedge[b=0\vee 3\cdot7^{b-1}\mid k]
  \wedge R(k).
\end{align*}
Here we used $\gcd(2,3^{a-1})=\gcd(2,3\cdot7^{b-1})=1$. The two right-hand sides differ only
in where the condition $2\mid k$ appears.
\begin{itemize}
  \item If $a,b\ge1$, each side contains the condition $2\mid k$.
  \item If exactly one of $a,b$ is positive, or both are $0$, then $r>1$ and $R(k)$ implies
    $2\mid k$.
\end{itemize}
Either way the two sides are equivalent for every $k$. Hence $\ord_m(x)=\ord_m(x+1)$, and
$f(m)\ge1$ after reducing $x$ modulo $m$. \qed

\begin{remark}
For $m=3^a7^b$ the orders are
$\ord_m(x)=\lcm(3^{a-1},6\cdot7^{b-1})$ and $\ord_m(x+1)=\lcm(2\cdot3^{a-1},3\cdot7^{b-1})$.
Both equal $2\cdot3^{\max(a-1,1)}\cdot7^{b-1}$. The construction was checked directly for every
odd $m\le 30000$.
\end{remark}

\section{Lean formalization}

The proof is formalized in Lean~4 (v4.33.1) with Mathlib, in about 520 lines; see the
accompanying file \texttt{A397683.lean}. The function $f$ is defined literally as in the OEIS.
Cohen's theorem is the explicit hypothesis
\begin{lstlisting}
def CohenConsecutive : Prop :=
  ∀ p : ℕ, p.Prime → 7 < p →
    ∃ y : ZMod p, orderOf y = p - 1 ∧ orderOf (y + 1) = p - 1
\end{lstlisting}
and the main theorem reads
\begin{lstlisting}
theorem conjecture (hC : CohenConsecutive) {m : ℕ}
    (hm : 1 < m) (hodd : Odd m) :
    f m = 0 ↔ (∃ a, m = 3 ^ a) ∨ (∃ b, m = 7 ^ b)
\end{lstlisting}
\texttt{\#print axioms} reports only \texttt{propext}, \texttt{Classical.choice} and
\texttt{Quot.sound}. In particular there is no \texttt{sorryAx}.

The lifting-the-exponent step uses Mathlib's \path{padicValNat.pow_sub_pow}. The lift
modulo $p^2$ uses \path{sq_dvd_add_pow_sub_sub}. The CRT gluing uses
\path{Nat.chineseRemainder} and \path{Nat.modEq_and_modEq_iff_modEq_mul}. The
reduction over the prime-power factorization uses \path{Nat.recOnPrimeCoprime}.

Cohen's theorem itself, which rests on character-sum estimates over finite fields, is not
formalized. To our knowledge it is not in Mathlib.

\begin{thebibliography}{9}
\bibitem{oeis} D.~S.~Gonz\'alez May, Sequence A397683, \emph{The On-Line Encyclopedia of
  Integer Sequences}, \url{https://oeis.org/A397683}, 2026.
\bibitem{cohen1} S.~D.~Cohen, Consecutive primitive roots in a finite field,
  \emph{Proc.\ Amer.\ Math.\ Soc.} 93 (1985), 189--197.
\bibitem{cohen2} S.~D.~Cohen, Consecutive primitive roots in a finite field~II,
  \emph{Proc.\ Amer.\ Math.\ Soc.} 94 (1985), 605--611.
\bibitem{cot} S.~D.~Cohen, T.~Oliveira e Silva and T.~Trudgian, On consecutive primitive
  elements in a finite field, \emph{Bull.\ London Math.\ Soc.} 47 (2015), 418--426;
  arXiv:1410.6210.
\bibitem{carella} N.~A.~Carella, Least consecutive pair of primitive roots,
  arXiv:2604.13090, 2026.
\bibitem{mathlib} The mathlib Community, The Lean mathematical library, \emph{CPP 2020}.
\end{thebibliography}

\end{document}
